\newpage
\section{为什么需要 tkz-elements？}

创建 \tkzNamePack{tkz-elements} 宏包，是为了解决在 TeX 中进行几何工作时的两个关键难题。第一，单靠 TeX 缺乏复杂几何计算所需的数值精度与灵活性。借助 \tkzname{Lua}，\tkzNamePack{tkz-elements} 为求交、投影、变换等运算提供了精确而稳定的计算。

第二，本宏包引入了面向对象的几何方法。点、向量、直线、圆、圆锥曲线等几何实体，都被建模为带有相关方法的结构化对象。这使得用户能够以更自然、更易读的方式来表达几何关系，贴近数学语言。

除了这两个核心思想之外，\tkzNamePack{tkz-elements} 还带来了以下几项额外的好处：

\begin{itemize}

 \item 简化计算：\tkzname{Lua} 的浮点运算与富有表现力的语法，使几何计算比在纯 TeX 中更易于书写和理解。

\item  可扩展性：得益于模块化设计和类的使用，新增几何对象与方法几乎毫不费力，使本宏包可扩展、可适应。

\item  互操作性：计算得到的点和数据能与 \TIKZ{} 或其他 TeX 绘图工具顺畅集成，让用户能够把计算与可视化有效地结合起来。

 \item 教学上的清晰：凭借其结构化且富有表现力的语法，\tkzNamePack{tkz-elements} 非常适合教学用途——在教学中，作图与记号的清晰性至关重要。
\end{itemize}

总而言之，\tkzNamePack{tkz-elements} 在 TeX 生态中提供了一个现代化的、面向对象的、由 Lua 驱动的几何作图框架。

\subsection{计算精度}
\subsubsection{\TIKZ{} 中的计算精度}
在 \TIKZ{} 中，表达式 \tkzimp{|veclen(x,y)|} 会计算 $\sqrt{x^2+y^2}$。
这一计算通过多项式逼近来实现，其思路借鉴了 \tkzimp{Rouben Rostamian} 的想法。

\pgfkeys{/pgf/number format/.cd,std,precision=5} \pgfmathparse{veclen(65,72)}
\begin{mybox}{}
\begin{verbatim}
   pgfmathparse{veclen(65,72)} \pgfmathresult
\end{verbatim}
\end{mybox}

 \tkzHand $\sqrt{65^2+72^2} \approx \pmpn{\pgfmathresult} $ 。

\subsubsection{Lua 中的计算精度}
可以如下定义一个 |luaveclen| 宏：

\begin{mybox}{}
\begin{verbatim}
\def\luaveclen#1#2{\directlua{tex.print(string.format(
'\percentchar.5f',math.sqrt((#1)*(#1)+(#2)*(#2))))}}
\end{verbatim}
\end{mybox}

以及

\begin{mybox}
\begin{verbatim}
\luaveclen{65}{72}
\end{verbatim}
\end{mybox}

得到：\tkzHand $\sqrt{65^2+72^2} = \pmpn{\luaveclen{65}{72}} $ {\color{red}！！}

\medskip
这个误差虽然在以百分之一磅为单位放置页面上的对象时微不足道，但对数学证明的结果来说却会成为问题。此外，这些不精确还会累积，并导致错误的作图。

\vspace{.5em}
为了解决这种精度不足的问题，我最初引入了 \tkzNamePack{fp}，随后又引入了 \tkzNamePack{xfp} 宏包。近来，随着 Lua\LATEX{} 的出现，我加入了一个 \tkzname{Lua} 选项，旨在用 \tkzname{Lua} 进行计算。

这是我创建这个宏包的主要动机，其次要目标是引入面向对象编程（OOP）并简化用 Lua 进行编程的方式。OOP 的概念促使我进一步探索它的各种可能性。

当时，我从 {\tkzimpbf{Nicolas Kisselhoff}} 那里收到了一些 Lua 编程示例，但起初我很难理解这些代码，于是我花时间耐心地学习 Lua。最终，我得以开发出 \tkzname{\tkznameofpack}，并融入了他的许多思想，我把它们改造后用于这个宏包。

\subsubsection{使用对象}
后来，我看到了 \tkzimp{Roberto Giacomelli}\footnote{\href{https://www.guitex.org/home/images/meeting2012/slides/presentazione_giacomell_guitmeeting_2012.pdf}{Grafica ad oggetti con LuaTEX}} 的一篇文章，内容是用 \tkzname{Lua} 和 \tkzNamePack{\TIKZ} 工具进行面向对象编程。这成为我灵感的第二个来源。这种方法不仅使编程可以一步一步地执行，而且对象的引入还在代码与几何之间建立了直接联系。结果，代码变得更加可读、明确，结构也更加良好。


\subsubsection{示例：阿波罗尼斯圆（新版本 2025/05/12）}
\label{ssub:example_apollonius_circle}

\begin{mybox}{问题：}
目标是找出一个与三角形的三个旁切圆都内切的圆。\end{mybox}


如需更多信息，可查阅 \href{https://mathworld.wolfram.com/ApolloniusCircle.html}{MathWorld} 上的相应条目。这个示例曾被用来测试 \tkzNamePack{tkz-euclide} 宏包。最初，用基本方法和现有工具得到的结果缺乏精度。得益于 \tkzNamePack{tkz-elements}，现在有了更强大、更精确的工具——而且它们也更易于使用。用 \tkzNamePack{tkz-euclide} 作图的核心原则保持不变：定义、计算、绘图和标注，都遵循一步一步的流程，与经典的尺规作图相呼应。这个版本利用了由 \tkzname{Lua} 所支持的最简便的作图方法。


\begin{mybox}
\begin{verbatim}
\directlua{
 init_elements()
 z.A = point(0, 0)
 z.B = point(6, 0)
 z.C = point(0.8, 4)
 T.ABC = triangle(z.A, z.B, z.C)
 z.N = T.ABC.eulercenter
 z.S = T.ABC.spiekercenter
 T.feuerbach = T.ABC:feuerbach()
 z.Ea, z.Eb, z.Ec = T.feuerbach:get()
 T.excentral = T.ABC:excentral()
 z.Ja, z.Jb, z.Jc = T.excentral:get()
 C.JaEa = circle(z.Ja, z.Ea)
% C.ortho = circle:radius(z.S, math.sqrt(C.JaEa:power(z.S)))
 C.ortho = C.JaEa:orthogonal_from(z.S) % 2025/05/12
 z.a = C.ortho.through
 C.euler = T.ABC:euler_circle()
 C.apo = C.ortho:inversion(C.euler)
 z.O = C.apo.center
 z.xa, z.xb, z.xc = C.ortho:inversion(z.Ea, z.Eb, z.Ec)}
\end{verbatim}
\end{mybox}

创建一个对象，就是定义它的属性和方法——也就是它的性质以及它所能执行的操作。对象一旦创建，就会通过存放在一个表中而与一个名字（或引用）相关联。这个表充当一个关联数组，把 \code{键}（引用）与 \code{值}（对象本身）联系起来。这些概念将在后面更详细地探讨。

例如，假设 \code{T} 是一个把三角形对象与键 \code{ABC} 关联起来的表。那么 \code{T.ABC} 就指向另一个包含该三角形属性的表——例如它的顶点、边或角——每个属性都通过一个特定的键来访问。这些属性在宏包内预定义好了，以支持几何运算。

\vspace{1em}
\begin{mybox}
\begin{verbatim}
 z.N = T.ABC.eulercenter
\end{verbatim}
\end{mybox}

\code{N} 是该点的名字，\code{eulercenter} 是三角形的一个属性。
\footnote{欧拉圆（即九点圆）的圆心，是每个三角形都具有的一个特征点。}

\begin{mybox}
\begin{verbatim}
 T.excentral = T.ABC:excentral()
\end{verbatim}
\end{mybox}

在这里，\tkzname{excentral} 是与 \tkzname{T.ABC} 对象相关联的一个方法。它定义了由三个旁切圆的圆心所构成的三角形。


\begin{tcolorbox}[
  title=已弃用代码,
  colframe=red!50,
  colback=red!5,
  coltitle=black,
  sharp corners,
  fonttitle=\bfseries,
  leftrule=1mm,
  rightrule=1mm,
  toprule=0.5mm,
  bottomrule=0.5mm
]
\texttt{注意：} 这段代码已被更优雅的版本取代。

有两行特别值得注意。第一行展示了 \tkzname{Lua} 卓越的精度如何使我们能够通过一次复杂的计算来定义半径。

根圆的半径由下式给出：
\[
\sqrt{\Pi(S,\mathcal{C}(Ja,Ea))}
\]
（即点 $S$ 关于以 |Ja| 为圆心、经过 |Ea| 的旁切圆的幂的平方根）。

\begin{mybox}{}
\begin{verbatim}
  C.ortho  = circle: radius (z.S,math.sqrt(C.JaEa: power(z.S)))
\end{verbatim}
\end{mybox}
\end{tcolorbox}

\begin{tcolorbox}[
  title=修订后的代码,
  colframe=cyan!60!black,
  colback=cyan!5,
  coltitle=white,
  sharp corners,
  fonttitle=\bfseries,
  leftrule=1mm,
  rightrule=1mm,
  toprule=0.5mm,
  bottomrule=0.5mm
]
\texttt{注意：} 以下代码以一种更优雅、更面向对象的方式取代了先前的版本。
先前的代码展示了可以手动执行的那类计算。然而，\tkzNamePack{tkz-elements} 为各种几何对象提供了丰富的方法。在与圆相关的方法中，有一个方法允许你在给定圆心的情况下，定义一个与另一个圆正交的圆。

我们希望得到一个与三角形三个旁切圆都正交的圆。它的圆心（$S$）就是这三个圆的根心。只需计算一个与其中之一正交的圆\footnote{给定一个主圆和两个辅圆，存在一个特定方法来定义与这三个圆都正交的根圆。}，以点 $S$ 为圆心即可。

第二行重要代码则是关于这个正交圆做一次反演。

\begin{mybox}{}
\begin{verbatim}
  C.ortho  = C.JaEa:orthogonal_from(z.S)
\end{verbatim}
\end{mybox}
\end{tcolorbox}

\vspace{1em}

最后，值得一提的是：把欧拉圆关于根圆做反演，就得到阿波罗尼斯圆。\footnote{九点圆，又称欧拉圆，与三个旁切圆都外切。切点构成费尔巴哈三角形。}

这个变换需要一个对象作为参数。该方法会自动检测对象类型（因为本宏包中所有对象都是有类型的），并据此选择合适的算法。

\begin{mybox}{}
\begin{verbatim}
  C.apo   = C.ortho:inversion(C.euler)
\end{verbatim}
\end{mybox}

现在所有相关的点都已定义好，接下来就该开始绘制几何路径了。为此，必须先创建相应的节点。这正是 \tkzMacro{tkz-elements}{tkzGetNodes} 宏的作用。

下一节专门讨论绘图，由 \tkzNamePack{tkz-euclide} 负责。

\begin{verbatim}
\begin{tikzpicture}
   \tkzGetNodes
   \tkzFillCircles[green!30](O,xa)
   \tkzFillCircles[teal!30](Ja,Ea Jb,Eb Jc,Ec)
   \tkzFillCircles[lightgray](S,a)
   \tkzFillCircles[green!30](N,Ea)
   \tkzDrawPoints(xa,xb,xc)
   \tkzDrawCircles(Ja,Ea Jb,Eb Jc,Ec S,a O,xa N,Ea)
   \tkzClipCircle(O,xa)
   \tkzDrawLines[add=3 and 3](A,B A,C B,C)
   \tkzDrawPoints(O,A,B,C,S,Ea,Eb,Ec,N)
   \tkzDrawSegments[dashed](S,xa S,xb S,xc)
   \tkzLabelPoints(O,N,A,B)
   \tkzLabelPoints[right](S,C)
\end{tikzpicture}
\end{verbatim}

\vspace{1em}
\directlua{
 init_elements()
 z.A = point(0, 0)
 z.B = point(6, 0)
 z.C = point(0.8, 4)
 T.ABC = triangle(z.A, z.B, z.C)
 z.N = T.ABC.eulercenter
 z.S = T.ABC.spiekercenter
 T.feuerbach = T.ABC:feuerbach()
 z.Ea, z.Eb, z.Ec = T.feuerbach:get()
 T.excentral = T.ABC:excentral()
 z.Ja, z.Jb, z.Jc = T.excentral:get()
 C.JaEa = circle(z.Ja, z.Ea)
 C.ortho = C.JaEa:orthogonal_from(z.S)
 z.a = C.ortho.through
 C.euler = T.ABC:euler_circle()
 C.apo = C.ortho:inversion(C.euler)
 z.O = C.apo.center
 z.xa, z.xb, z.xc = C.ortho:inversion(z.Ea, z.Eb, z.Ec)
}
\begin{minipage}{\textwidth}

\vspace{1em}
\begin{center}
\begin{tikzpicture}[ scale  = .6]
   \tkzGetNodes
   \tkzFillCircles[purple!50](O,xa)
   \tkzFillCircles[orange!50](Ja,Ea Jb,Eb Jc,Ec)
   \tkzFillCircles[yellow](S,a)
   \tkzFillCircles[purple!30](N,Ea)
   \tkzDrawPoints(xa,xb,xc)
   \tkzDrawCircles(Ja,Ea Jb,Eb Jc,Ec S,a O,xa N,Ea)
   \tkzClipCircle(O,xa)
   \tkzDrawLines[add=3 and 3](A,B A,C B,C)
   \tkzDrawPoints[red](O,A,B,C,S,Ea,Eb,Ec,N)
   \tkzDrawSegments[dashed](S,xa S,xb S,xc)
   \tkzLabelPoints(O,N,A,B)
   \tkzLabelPoints[right](S,C)
\end{tikzpicture}
\end{center}
\end{minipage}


\endinput
